\label{ch:conclusion}
% Chapter role anchor and validation scope recap
This chapter summarizes the thesis contribution, situates it in the prior-work design space, and identifies open problems and future work. This thesis investigates simultaneous \gls{emg} and \gls{dc} \gls{eda} acquisition from the same muscle site. Co-locating these modalities introduces a ground loop between the two measurement circuits. The custom front-end places the \gls{emg} and \gls{eda} chains in separate ground domains, with reinforced galvanic isolation between each chain and a common data acquisition unit. The system was validated on a single participant at the palm across five recording sessions. The biceps clinical site is not directly validated.

% Three central findings: signal chain works, cross-modality asymmetric, pattern reproduces
Three central findings emerge from the validation. Each front-end individually resolves its target signal against the post-isolation noise floor. The two cross-modality directions behave differently. \gls{eda} operation produces no detectable elevation in the \gls{emg} amplitude metrics, so combined acquisition is technically supported in the \gls{eda} to \gls{emg} direction. \gls{emg} operation leaves a measurable signature on the \gls{eda} signal. A reproducible downward dip in the \gls{eda} signal follows each \gls{emg} cue, and its mechanism remains unidentified. The directional pattern reproduces across sessions, with one session inflating between-session variability.

% Contribution position, originality, design-space partition, mechanism, validation confirmation, clinical significance
The contribution sits in a previously unaddressed region of the multimodal design space. No prior work co-locates \gls{emg} and \gls{dc} \gls{eda} at the same muscle site with galvanic isolation. The thesis adds one point to a design space organized along three dimensions: modality combinations, excitation regimes, and anatomical sites. The \gls{dc} excitation current returns through the \gls{emg} ground electrode when both chains share a ground domain, and galvanic isolation breaks the shared-ground paths. The validation confirms the technical precondition for combined acquisition. The design is a step toward multimodal stress and pain monitoring in non-verbal patient populations.

% Limitations, future-work directions, headline reaffirmation, forward closer
The contribution has three limitations: the \gls{emg} to \gls{eda} coupling mechanism is unidentified, validation is at the palm only, and the cross-session sample size is $n = \num{3}$. Section~\ref{sec:disc-future} addresses each: bench characterization of the isolation stage and cross-modality coupling, validation refinements through a biceps recording campaign, and generalization across the prior-work design space. Galvanic isolation supports co-located \gls{emg} and \gls{dc} \gls{eda} acquisition at the same muscle site. Future work extends this foundation along clinical and design-space dimensions.
\begin{comment}
    
\begin{itemize}
    \item Summary of you work
    \item Your work in context and a broader perspective
    \begin{itemize}
        \item What is the main contribution of the thesis on the problem
        \item Why is the contribution original?
        \item Why is the contribution non-trivial?
    \end{itemize}
    \item Reflection on the topic and process
    \begin{itemize}
        \item Significance and limitations
        \item What went wrong?
        \item Identify open problems
        \item Point to future research possibilities 
    \end{itemize}
\end{itemize}
\end{comment}